\chapter{Derivations for Decision under Uncertainty}
\label{annex:p17-rationality-8-4}

\begin{annexitem}{annex:p17.rat84-criteria}{thm:p17-rat84-criteria}{The source's supplied-prior comparisons and the separating menu}
\begin{proof}
The expected utility of \(1_E0\) is \(p(E)\). On
\(\mathcal C^*\), both unknown colors have lower probability
\(.3\) and upper probability \(.7\); both known colors have
probability \(.5\). All endpoints are attainable: for example
independent known/unknown colors with either extreme unknown
proportion define admissible joint vectors in the four-state simplex.
Thus the maxmin values are exactly those stated, and the alpha value
for an unknown-color bet is
\(\alpha(.3)+(1-\alpha)(.7)=.7-.4\alpha\).
Comparing it with \(.5\) gives strict known preference iff
\(\alpha>1/2\). At unknown probability \(.3\) the known bet
is better; at \(.7\) it is worse. Neither unanimous orientation
therefore holds. The same argument applies to the complementary color.

For the three-act menu write a two-state prior as \((t,1-t)\).
Its expected utilities are \(t,1-t,.4\). Every pair has opposite
rankings at some priors, so no act dominates another under unanimity.
The first act is best for \(t\geq.6\), and the second for
\(t\leq.4\). The constant act can never be best because
\(\max(t,1-t)\geq.5>.4\). Thus precisely the first two are
E-admissible. Their worst values are both zero, whereas the
constant's worst value is \(.4\). Maxmin uniquely selects it.
Their alpha values are \(1-\alpha,1-\alpha,.4\), proving
the threshold \(.6\) and all tie cases. The contradiction
between criteria uses their actual distinct choice definitions,
not a different prior set or utility for each rule.
\end{proof}
\end{annexitem}

\begin{annexitem}{annex:p17.rat84-axioms}{thm:p17-rat84-axioms}{Stake independence, hedging and the continuation counterexample}
\begin{proof}
For a bet \(x_Ey\), with \(x>y\), its expected utility is
\(y+(x-y)p(E)\). Its maxmin value is therefore
\(y+(x-y)\min_{\mathcal C}p(E)\), and its fixed-alpha value
is \(y+(x-y)[\alpha\min p(E)+(1-\alpha)\max p(E)]\).
Comparing two events with the same stakes cancels \(y\) and
the positive factor \(x-y\); the ranking is independent of
the prizes. This proves P4 in both models.
For unanimity, \(EU_p(f_Eh)-EU_p(g_Eh)\) contains only terms
on \(E\), independently of the common outside act \(h\).
Its sign at every retained prior is consequently unchanged by
replacing that outside act, proving STP even when some pairs
remain incomparable.

For utility vectors \(v,w\), every retained prior gives
\[
 EU_p(\lambda v+(1-\lambda)w)
 \geq\lambda\min_q EU_qv+(1-\lambda)\min_q EU_qw.
\]
Minimize the left side to prove maxmin concavity. If the two
disjoint event bets have the same maxmin value, their half
utility mixture is the source's hedge and has at least that
value. With two point priors, the complementary bets each
have minimum zero and maximum one. Their alpha values are
\(1-\alpha\); their hedge is constant \(.5\). When
\(\alpha<1/2\), this is strictly worse, disproving UA.

For the dynamic example, each original prior has expected
utility \(.5\) for \(f=(1,0,0,1)\), so its ex ante maxmin
value is \(.5>.4\). Both priors give probability \(.5\)
to \(E\) and to \(E^c\), so all conditional priors exist.
On \(E\) their conditionals are the two point masses on
states 1 and 2; on \(E^c\) they are the point masses on
states 3 and 4. In either conditional problem the act's
minimum utility is zero, while the constant act remains
\(.4\). Hence both continuations reverse the ex ante
comparison under the stated updating procedure. The original
scalar ranking is nevertheless complete and transitive.
This explicitly distinguishes a static property from the
additional requirement of dynamic consistency.
\end{proof}
\end{annexitem}

\begin{annexitem}{annex:p17.rat84-comparative}{thm:p17-rat84-comparative}{Both directions of the finite prior-inclusion characterization}
\begin{proof}
Write \(V_i(v)=\min_{p\in\mathcal C_i}p\cdot v\).
The behavioral condition is equivalent to \(V_1(v)\leq V_2(v)\)
for every utility vector \(v\). That inequality implies the
condition directly. Conversely, if \(V_1(v)>V_2(v)\), choose
an available sure utility strictly between the two values;
agent 1 accepts \(v\) over it while agent 2 does not.
Thus the domain of sure consequences is essential to the converse.

If \(\mathcal C_2\subseteq\mathcal C_1\), a minimum over
the larger set is at most the minimum over the smaller set,
giving the required value inequalities. Suppose instead
\(p_2\in\mathcal C_2\setminus\mathcal C_1\). By compactness
choose a point \(q\in\mathcal C_1\) minimizing its Euclidean
distance from \(p_2\); its distance is positive since the
set is closed. For any \(p\in\mathcal C_1\), convexity
keeps \(q+t(p-q)\) in that set for \(0\leq t\leq1\).
Minimality of the squared distance at \(t=0\) gives
\((q-p_2)\cdot(p-q)\geq0\), by expanding the quadratic
and dividing by positive \(t\) before letting it decrease
to zero. With \(v=q-p_2\), consequently
\[
 p\cdot v\geq q\cdot v>p_2\cdot v
             \quad\text{for every }p\in\mathcal C_1.
\]
The strict gap is \(\|q-p_2\|^2\). Therefore
\(V_1(v)>p_2\cdot v\geq V_2(v)\), contradicting the
behavioral condition. This proves inclusion.

A linear expected utility has the same extrema on a set,
its convex hull, and its closure, whenever these extrema
are taken on the compact closed hull. Thus canonical hulls,
not arbitrary raw lists, are the sets identified by this
argument. The separating vector need not be among a few
elicited bets; that explains the rich-domain qualification.
\end{proof}
\end{annexitem}

\begin{annexitem}{annex:p17.rat84-choquet}{thm:p17-rat84-choquet}{Signed Choquet values, the greedy core and the precise naive countercase}
\begin{proof}
The formula is unchanged by adding a common constant except
for adding that constant to the baseline. For a nonnegative
vector it equals the area under the finite step function
\(t\mapsto\nu(\{s:v(s)\geq t\})\), \(t\geq0\).
No integration theorem is needed: partition the positive
axis at the finitely many utility values and sum rectangle
areas. If \(v\leq w\), every upper-level set for \(v\)
is included in the corresponding one for \(w\), so each
step height is no larger. For signed vectors, add the same
large constant to both and cancel it. This proves monotonicity.
Sorting all individual state values, with repetitions, gives
the same formula: each repeated-value increment is zero.
The order of tied states therefore has no effect.
For a two-valued bet the formula is exactly
\(y+(x-y)\nu(E)\), proving the P4 statement.

To prove the core representation, order all states as
\(s_1,\ldots,s_n\) by nonincreasing utility and put
\(A_i=\{s_1,\ldots,s_i\}\), \(A_0=\varnothing\).
Define the greedy probability
\[
               p^*(s_i)=\nu(A_i)-\nu(A_{i-1}).
\]
It is nonnegative by monotonicity and sums to one by
normalization. Convexity gives increasing marginal increments:
if \(B\subseteq A\) and \(s\notin A\), then
\[
 \nu(A\cup\{s\})-\nu(A)
                  \geq\nu(B\cup\{s\})-\nu(B).
\]
This is the convex-capacity inequality applied to
\(A\) and \(B\cup\{s\}\), whose union is
\(A\cup\{s\}\) and intersection is \(B\).
For any \(E\subseteq S\), apply it at each \(s_i\in E\)
with \(A=A_{i-1}\), \(B=E\cap A_{i-1}\). Summing the
resulting inequalities telescopes to
\[
 p^*(E)\geq\sum_{s_i\in E}
  [\nu(E\cap A_i)-\nu(E\cap A_{i-1})]=\nu(E).
\]
Thus \(p^*\) is in the core, proving nonemptiness.
For every probability \(p\), summation by parts gives
\[
 p\cdot v=v(s_n)+\sum_{i=1}^{n-1}
                    [v(s_i)-v(s_{i+1})]p(A_i).
\]
If \(p\) is in the core, all nonnegative-increment terms
are at least the corresponding Choquet terms. For \(p^*\),
\(p^*(A_i)=\nu(A_i)\) for every \(i\), so equality holds.
It attains the minimum and completes the core identity.
The earlier maxmin-concavity proof then supplies UA.

For the nonconvex example, each one-state bet has Choquet
value \(.75\), while the constant half act has value \(.5\).
The capacity fails convexity because
\(\nu(S)+\nu(\varnothing)=1<.75+.75\). For the naive
outcome-event substitution with singleton values \(.1,.1\),
the constant act has one outcome event \(S\) and value
\(1\nu(S)=1\); the improved act has two different
outcome events and value \(1.1(.1)+1.2(.1)=.23\).
Its Choquet value is \(1.1+(1.2-1.1)(.1)=1.11\).
On a fixed state partition, in contrast, the naive coefficients
are nonnegative, so that sum is monotone; for the constant
one it gives \(.1+.1=.2\), not one. This is a different
defect and a different substitution operation.
\end{proof}
\end{annexitem}

\begin{annexitem}{annex:p17.rat84-confidence}{thm:p17-rat84-confidence}{The source's nested intervals and an unrestricted-selector limit}
\begin{proof}
At the point prior \(p_0\), both black bets have expectation
\(.5\). At the large high-confidence set the known marginal
is still \(.5\), while the least unknown-black marginal
is \(.1\). Scaling the winning utility to 100 therefore
gives 50 and 10. The same beliefs family has produced
indifference at one stake and a strict preference at another,
so P4's stake-invariant comparison fails in this procedure.

For the independent boundary example, minimization of
\((2,1)\) over all two-state priors gives 1, and evaluation
of \((1.9,.9)\) at its prescribed equal prior gives 1.4.
Nevertheless \(2>1.9\) and \(1>.9\). The reversal arises
from changing the selected set, not from a failure of
ordinary expectations to respect dominance. No axiom
restricting \(D\) was imposed in this counterexample,
so it does not ascribe the reversal to all confidence theories.
\end{proof}
\end{annexitem}

\begin{annexitem}{annex:p17.rat84-cardinal}{thm:p17-rat84-cardinal}{Explicit penalty and reciprocal calculations and existence}
\begin{proof}
On the two-prior candidate list the penalty values for
utilities \((0,1)\), constant \(.3\), \((.4,.6)\), and
\((.1,.6)\) are respectively
\[
 \min(.2,1)=.2,\quad\min(.5,.3)=.3,\quad
 \min(.6,.6)=.6,\quad\min(.3,.6)=.3.
\]
The reciprocal values under \((.5,1)\) are respectively
\[
 \min(0,1)=0,\quad\min(.6,.3)=.3,\quad
 \min(.8,.6)=.6,\quad\min(.2,.6)=.2.
\]
The raw maxmin of \((.4,.6)\) is \(.4\). The equality of
one pair of values therefore does not equate the rules.

In the general finite-state case \(p\mapsto EU_p f\)
is continuous and bounded on compact \(\Delta\). Adding
a nonnegative lower-semicontinuous penalty gives a
lower-semicontinuous objective bounded below. It has a
finite feasible value at a prior where the penalty is zero.
A minimizing sequence has a convergent subsequence by
compactness; lower semicontinuity at its limit attains
the infimum. For constant utility \(z\), the value is
\(z+\min c=z\). Taking the explicitly extended indicator
penalty restricts the finite objective exactly to the closed
nonempty \(\mathcal C\), giving maxmin there.

With positive continuous \(\varphi\) on compact
\(\Delta\), its minimum is positive and the ratio is
continuous, so the minimum exists. If \(\max\varphi=1\)
and \(z\geq0\) is constant, \(z/\varphi(p)\geq z\)
and equality is attainable, preserving that sure utility.
For \(z<0\), the ratio becomes smaller as
\(\varphi\) decreases and in general its minimum is
\(z/\min\varphi\), not \(z\). If zeros are allowed,
neither division nor the signed comparison is defined
by the bare expression without a further convention.
These are the precise boundary conditions of the local
reciprocal calculation.
\end{proof}
\end{annexitem}

\begin{annexitem}{annex:p17.rat84-reduction}{thm:p17-rat84-reduction}{Second-order averaging and the nonaffine criterion}
\begin{proof}
Finite sums commute, giving
\[
 \sum_j\mu_jEU_{p_j}f
 =\sum_s\left(\sum_j\mu_jp_j(s)\right)v_f(s)
 =EU_{\bar p}f.
\]
The mixture \(\bar p\) is a normalized nonnegative
probability vector. If \(p(E)=p(F)\), the common weighting
function gives \(w(p(E))=w(p(F))\); the equal-prize
Choquet bet formula then gives identical values. This
proves both reduction limitations.

For the displayed smooth case, known-urn expected utility
is \(.5\) at both priors. Unknown-urn expected utility
is 0 at one and 1 at the other. The complementary
color swaps these two values, leaving their average
unchanged. Under the square root, the known value is
\(\sqrt{.5}\), strictly greater than
\((\sqrt0+\sqrt1)/2=.5\). Under the square it is
\(.25<(.0^2+1^2)/2=.5\). An affine transform has
\(\sum_j\mu_j\varphi(EU_{p_j}f)
=\varphi(\sum_j\mu_jEU_{p_j}f)\), so it depends only
on the reduced expected utility and the equal means tie.

For finitely many utility values and positive weights,
concavity gives
\(\sum_j\mu_j\varphi(z_j)\leq
\varphi(\sum_j\mu_jz_j)\), by successively grouping
two weighted values and applying the defining concavity
inequality. With a strictly concave transform and at
least two distinct values of positive weight, one of
those groupings is strict. A constant expected utility
with the same mean therefore receives the strictly
higher transformed value. The result concerns the
given second-order model; it does not infer \(\mu\)
or the transform from the first-order mean alone.
\end{proof}
\end{annexitem}
